\section{Variational Principle and Path-Independent Convergence to Action-Cells}
\label{sec:variational-principle}

%=============================================================================
\subsection{Muscle Action as a Variational Problem}
\label{subsec:variational-formulation}

The execution of a motor action (e.g., lifting a weight, pressing a button, reaching to a target) can be formalized as an extremization problem over the space of motor commands. Instead of executing a pre-specified motor program, the nervous system navigates a trajectory through motor command space such that the **action functional** is minimized.

\begin{definition}[Motor Action Functional]
\label{def:action-functional}
For a motor command trajectory $\mathbf{u}(t) = (u_{\text{spinal}}(t), u_{\text{cortical}}(t), u_{\text{cerebellar}}(t), \ldots)$ from $t=0$ to $t=T$, the motor action functional is:
\begin{equation}
  \mathcal{S}[\mathbf{u}] = \int_0^T \mathcal{L}(\mathbf{u}, \dot{\mathbf{u}}, \mathbf{x}, \dot{\mathbf{x}}) \, dt,
\end{equation}
where:
\begin{itemize}[noitemsep]
\item $\mathcal{L}$ is the motor Lagrangian (energy cost functional)
\item $\mathbf{x}(t) = (\text{muscle force}, \text{joint angle}, \text{limb velocity})$ is the motor state
\item $\dot{\mathbf{x}}(t)$ is the rate of change of motor state
\end{itemize}

A motor action trajectory is one that satisfies:
\begin{equation}
  \delta \mathcal{S}[\mathbf{u}] = 0,
\end{equation}
i.e., the action is stationary (minimized or maximized) with respect to variations in the motor command.
\end{definition}

\begin{principle}[Minimality of Motor Action]
The nervous system selects motor command trajectories that minimize (or near-minimize) the action functional $\mathcal{S}[\mathbf{u}]$. This is not a conscious optimization; it emerges from the physical constraints of the motor system (ATP limitation, mechanical efficiency, proprioceptive feedback).
\end{principle}

%=============================================================================
\subsection{The Muscle Lagrangian from First Principles}
\label{subsec:muscle-lagrangian}

We now derive the motor Lagrangian for muscle contraction, incorporating:
1. ATP consumption cost (oxygen-dependent)
2. Mechanical work done on the skeleton
3. Proprioceptive error penalty
4. Synaptic state constraints (residual neurotransmitter)

\begin{definition}[The Muscle Lagrangian]
\label{def:muscle-lagrangian}
The motor Lagrangian for a single muscle or muscle group is:
\begin{equation}
\boxed{
\mathcal{L} = \underbrace{\frac{1}{2}m\dot{x}^2}_{\text{kinetic}} - \underbrace{F_{\text{ext}} x}_{\text{work}} - \underbrace{\alpha_{\text{ATP}} \frac{P_{\text{ATP}}(t)}{P_{\text{ATP,max}}} F_{\text{muscle}}}_{\text{ATP cost}} - \underbrace{\beta \|[\text{prop}]_{\text{actual}} - [\text{prop}]_{\text{expected}}\|^2}_{\text{proprioceptive error}} + \underbrace{\gamma \sigma_{\text{residue}}}_{\text{synaptic residue}}
}
\label{eq:muscle-lagrangian}
\end{equation}

where:
\begin{itemize}[noitemsep]
\item $m$ is the effective inertial mass (limb + muscle)
\item $\dot{x}$ is limb velocity
\item $F_{\text{ext}}$ is the external force (load being lifted)
\item $x$ is limb displacement
\item $P_{\text{ATP}}(t)$ is the ATP production rate, $P_{\text{ATP,max}}$ is maximal rate (determined by oxygen availability per Section~\ref{sec:cellular-architecture})
\item $F_{\text{muscle}}$ is the muscle force generated
\item $[\text{prop}]_{\text{actual}}$ is the sensory proprioceptive signal (muscle spindle, tendon organ, joint receptor)
\item $[\text{prop}]_{\text{expected}}$ is the internally predicted proprioceptive signal (efference copy from motor cortex)
\item $\sigma_{\text{residue}}$ is the synaptic residual neurotransmitter concentration (from Section~\ref{subsec:residue-anchor})
\end{itemize}
\end{definition}

\begin{theorem}[ATP Cost Scales with Oxygen Availability]
\label{thm:atp-oxygen-coupling}
The ATP production rate $P_{\text{ATP}}(t)$ is determined by local oxygen concentration:
\begin{equation}
  P_{\text{ATP}}(t) = P_{\text{max}} \cdot \frac{[\text{O}_2](t)}{K_{m,\text{O}_2} + [\text{O}_2](t)},
\end{equation}

where $K_{m,\text{O}_2} \sim 1$--$10$ $\mu$M is the oxygen Michaelis constant for oxidative phosphorylation. Therefore, regions of muscle operating at low oxygen are energetically expensive (high $\alpha_{\text{ATP}} F_{\text{muscle}} / P_{\text{ATP}}$), creating a geometric preference for oxygen-rich (superficial) pathways.
\end{theorem}

\begin{proof}
By Theorem~\ref{thm:oxygen-hierarchy} (Section~\ref{subsec:oxygen-timekeeper}), enzyme rates are governed by Michaelis-Menten kinetics with oxygen as substrate. The ATP synthase in mitochondria operates at rates proportional to local oxygen availability. Deep muscle fibers, being oxygen-poor, cannot generate ATP fast enough to sustain high-force contractions, so the motor system gravitates toward recruiting oxygen-rich (superficial, Type I) fibers first. $\square$
\end{proof}

\begin{corollary}[Energy Cost Encodes Spatial Oxygen Distribution]
The Lagrangian's ATP term $-\alpha_{\text{ATP}} P_{\text{ATP}} F$ is not uniform in space. Regions of high oxygen concentration (near capillaries) appear as deep potential wells in the Lagrangian; regions of low oxygen appear as barriers. The motor system navigates toward low-ATP-cost actions by default, which means preferring oxygen-rich pathways.
\end{corollary}

%=============================================================================
\subsection{Action-Cells: Behavioral Goals in Motor State Space}
\label{subsec:action-cells-motor}

A motor action-cell is a region of motor state space where a specific behavioral goal is achieved. Unlike the intracellular action-cells of Section~\ref{subsec:intracellular-circulation}, motor action-cells are defined behaviorally.

\begin{definition}[Motor Action-Cell]
\label{def:motor-action-cell}
A motor action-cell for a task $\tau$ is a region $\mathcal{A}_\tau \subset \text{Motor State Space}$ characterized by:
\begin{enumerate}[label=(\roman*), noitemsep]
\item A target force range: $F_{\tau,\min} \leq F_{\text{muscle}} \leq F_{\tau,\max}$ (e.g., grip strength for object lifting)
\item A target velocity range: $v_{\tau,\min} \leq |\dot{x}| \leq v_{\tau,\max}$ (e.g., reaching speed)
\item A target endpoint region: the limb endpoint is within distance $\delta_x$ of the target location
\item A proprioceptive verification condition: sensory feedback matches expected feedback within tolerance $\delta_{\text{prop}}$
\end{enumerate}

All motor commands that drive the system into $\mathcal{A}_\tau$ result in successful task completion.
\end{definition}

\begin{example}[Action-Cell for Lifting a Mug]
\label{ex:mug-lift}
The task is to lift a coffee mug (mass $\approx 0.25$ kg) from a table to the mouth.

The action-cell $\mathcal{A}_{\text{lift}}$ is defined by:
\begin{itemize}[noitemsep]
\item Force: $F_{\text{grip}} \in [2, 5]$ N (enough to support mug without crushing)
\item Vertical velocity: $v_z \in [0.05, 0.5]$ m/s (slow enough to not spill, fast enough to be efficient)
\item Endpoint: hand position within $\pm 2$ cm of mouth
\item Proprioception: predicted elbow angle matches actual elbow angle within $\pm 5°$
\end{itemize}

Motor commands satisfying these constraints all achieve successful lifting, even though they may have arrived via different neural pathways, different muscle-recruitment patterns, and different proprioceptive feedback signatures.
\end{example}

%=============================================================================
\subsection{Path-Independent Convergence to Motor Action-Cells}
\label{subsec:convergence-motor}

We now apply the path-independent convergence theorem (Theorem~\ref{thm:convergence} from asymptotic-categorical-junctions) to motor control.

\begin{theorem}[Path-Independent Convergence to Behavioral Action-Cells]
\label{thm:convergence-motor}
Let $\mathbf{u}_1(t), \mathbf{u}_2(t), \ldots, \mathbf{u}_n(t)$ be motor command trajectories originating from distinct neural circuits or feedback conditions. If all trajectories satisfy the motor sufficiency condition:
\begin{equation}
  \exists\, t^* : \|\nabla \mathcal{L}(\mathbf{u}(t))\| < \varepsilon_{\text{motor}} \quad \forall\, t > t^*,
\end{equation}

then all trajectories converge to the same motor action-cell $\mathcal{A}_\tau$:
\begin{equation}
  \mathbf{x}_i(t \to T) \in \mathcal{A}_\tau \quad \forall\, i \in \{1, \ldots, n\}.
\end{equation}

That is, multiple different motor commands generate identical behavioral outcomes.
\end{theorem}

\begin{proof}
By the motor sufficiency condition, all trajectories enter the basin of attraction $B_{\varepsilon}(\mathcal{A}_\tau)$ of the action-cell at time $t^*$. Within this basin, the Lagrangian $\mathcal{L}$ is minimized at the action-cell boundary. The gradient flow $\dot{\mathbf{x}} = -\nabla \mathcal{L}$ drives the system monotonically toward lower values of $\mathcal{L}$. By compactness of the motor state space (bounded limb positions, bounded muscle forces) and continuity of $\mathcal{L}$, all trajectories entering the basin converge to the action-cell. $\square$
\end{proof}

\begin{corollary}[Multiple Pathways Converge to the Same Action]
\label{cor:multiple-pathways}
Consider three motor command pathways to achieve the same action (e.g., lifting the mug):

\begin{itemize}[noitemsep]
\item \textbf{Pathway A} (Spinal Reflex): Load-sensitive feedback from muscle spindle directly drives motor neurons. Fast, low-latency, but crude.
\item \textbf{Pathway B} (Corticomotor): Motor cortex issues descending commands based on visual target. Slower, but precise target-reaching.
\item \textbf{Pathway C} (Cerebellar Correction): Cerebellum compares expected (efference copy) with actual (sensory) proprioception and issues error-correction signals.
\end{itemize}

All three pathways, if they satisfy the motor sufficiency condition, drive muscle contraction into the action-cell $\mathcal{A}_{\text{lift}}$ defined in Example~\ref{ex:mug-lift}. The task outcome (mug successfully lifted) is identical despite the three distinct neural circuits and distinct motor command trajectories.

This is the essence of motor redundancy and robustness: the behavior is over-specified (too many control pathways), but they all converge to the same outcome.
\end{corollary}

%=============================================================================
\subsection{Sufficiency Principle: Details Don't Matter, Goals Do}
\label{subsec:sufficiency-principle}

The convergence theorem implies a profound principle for motor control and, by extension, for cognition:

\begin{principle}[Sufficiency Principle in Motor Control]
\label{prin:sufficiency-motor}
A motor command need not be ``correct'' in an absolute sense. It need only satisfy the sufficiency condition: the gradient of the Lagrangian is small enough that the system is drawn into the appropriate action-cell.

Implications:
\begin{enumerate}[label=(\arabic*), noitemsep]
\item \textbf{Graceful Degradation}: Sensory information can be noisy, delayed, or incomplete. As long as it's sufficient to keep the system on-course toward the action-cell, the behavior succeeds.
\item \textbf{Internal Model Inaccuracy}: The motor cortex's internal model of the body (the predicted proprioception) need not be accurate. It need only be sufficient to drive the system toward the goal.
\item \textbf{Pathway Flexibility}: Different spinal circuits, different cerebellar gain settings, different proprioceptive sensitivities can all work, provided they satisfy sufficiency.
\item \textbf{No Perfect Execution}: Evolution did not optimize for perfect, novel movements each time. It optimized for sufficient, good-enough movements that work under realistic conditions (noise, fatigue, variability).
\end{enumerate}
\end{principle}

%=============================================================================
\subsection{Poincaré Deviation in Muscle: No Two Contractions Are Identical}
\label{subsec:poincare-muscle}

By Theorem~\ref{thm:poincare} (Poincaré Deviation, from asymptotic-categorical-junctions), feedback loops in the motor system do not exhibit exact recurrence. Each muscle contraction is unique.

\begin{theorem}[Poincaré Deviation in Motor Loops]
\label{thm:poincare-motor}
Consider a repeated motor action: lifting the same mug from the same location, executed at times $t_1, t_2, \ldots, t_n$. The motor command trajectory $\mathbf{u}_i(t)$ at time $t_i$ and the trajectory $\mathbf{u}_j(t)$ at time $t_j$ ($j > i$) satisfy:
\begin{equation}
  \mathbf{u}_i \neq \mathbf{u}_j \quad \text{with strictly positive deviation} \quad \delta = \min_t d_{\text{motor}}(\mathbf{u}_i(t), \mathbf{u}_j(t)) > 0.
\end{equation}

This is not a failure of motor control; it is a feature. Multiple re-executions generate distinct trajectories, which encode the history of motor action.
\end{theorem}

\begin{proof}
By Theorem~\ref{thm:poincare}, exact recurrence of a feedback loop requires the motor command to return exactly to its prior state. This requires:
\begin{enumerate}[noitemsep]
\item Synaptic state identical (same residual NT, same synaptic weights)
\item Proprioceptive state identical (same muscle spindle firing, same Golgi tendon organ signal)
\item Oxygen distribution identical (impossible, as metabolism has changed)
\item Cerebellar weight state identical (learning would have occurred)
\end{enumerate}

For measure-theoretic reasons (the set of exactly-recurrent trajectories has measure zero), repetition of a motor action generates a new motor command trajectory. Each execution is a new path. $\square$
\end{proof}

\begin{corollary}[Motor Memory is Trajectory History, Not Repetition]
\label{cor:motor-memory}
The ``memory'' of a movement is not a stored motor program. It is the trajectory taken to execute that movement. When you lift the mug again tomorrow, your motor cortex does not retrieve a stored program; it navigates a new trajectory that resembles the prior trajectory (because the physical constraints, the limb mechanics, and the learned cerebellar weights have shaped the motor landscape), but is never identical.

This explains:
\begin{itemize}[noitemsep]
\item \textbf{Motor Improvement with Practice}: Practice refines the motor landscape (cerebellar weights, motor cortical gain fields) so that each subsequent trajectory is more efficient, more accurate.
\item \textbf{Trial-to-Trial Variability}: Even under constant conditions, repeated movements vary because synaptic states differ (residue changes, neuromodulators fluctuate).
\item \textbf{Unique Movement Signatures}: Each individual's movement pattern is unique (fingerprint-like), because the trajectory history is unique.
\end{itemize}
\end{corollary}

%=============================================================================
\subsection{Multiple Thoughts, Single Action: Convergence at the Behavioral Terminus}
\label{subsec:multiple-thoughts}

The motor convergence theorem can be restated at a higher cognitive level:

\begin{theorem}[Thought-Action Convergence]
\label{thm:thought-convergence}
Consider a single behavioral goal (e.g., ``lift the mug''). Multiple distinct thought patterns and neural circuit states can all generate the same motor action:

\begin{itemize}[noitemsep]
\item \textbf{Thought A}: Visual, goal-directed thought (``I see the mug, I want it'')
\item \textbf{Thought B}: Habitual, procedural thought (``I've lifted mugs 1000 times, this is automatic'')
\item \textbf{Thought C}: Emotionally-driven thought (``I'm thirsty, I need coffee'')
\item \textbf{Thought D}: Sensorimotor thought (``My proprioception says my arm is ready to grasp'')
\end{itemize}

All four thought patterns (distinct trajectories through cerebral cortex, hippocampus, amygdala, and motor systems) converge to the identical motor action: arm extension, grip force, lift. The thought pattern is irrelevant; only the motor outcome matters.
\end{theorem}

\begin{proof}
By Theorem~\ref{thm:convergence-motor}, all motor command trajectories that satisfy the motor sufficiency condition converge to the action-cell. Different thought patterns generate different motor command trajectories (different neural circuits, different cortical sources), but if all satisfy sufficiency, they all reach the same behavioral outcome. The thought is the internal trajectory; the action is the terminal point. Multiple trajectories can end at the same terminus. $\square$
\end{proof}

\begin{corollary}[Consciousness as Convergence Verification]
\label{cor:consciousness-convergence}
Consciousness emerges at the point where multiple distinct internal model trajectories (thoughts, sensations, emotional states) are verified to converge on the same sensory outcome (action result verified by feedback).

When you think about lifting the mug:
- Motor cortex predicts the proprioceptive outcome
- Cerebellum estimates the trajectory
- Proprioceptive system reports actual trajectory
- Visual system confirms mug location changed

Consciousness is the intersection of these multiple convergent trajectories being **verified** against reality (external feedback). If the predicted and actual outcomes diverge, consciousness flags an error and the system corrects.

In dreams, there is no external verification. Multiple motor command trajectories are generated (you attempt to run, to fly, to scream), but with no sensory feedback to verify outcomes, they diverge and rebound chaotically. No true convergence, hence no consciousness.
\end{corollary}

%=============================================================================
\subsection{The Sufficiency Principle Unifies All Regimes}
\label{subsec:sufficiency-unification}

We now state the central unifying principle:

\begin{principle}[Universal Sufficiency]
\label{prin:universal-sufficiency}
All dynamical regimes discussed in this paper---from oxygen-driven enzyme kinetics (Section~\ref{sec:cellular-architecture}) through motor neuron recruitment (Section~\ref{subsec:motor-recruitment}) to behavioral action execution (current section)---obey the same principle:

\textbf{Multiple distinct trajectories through a system's state space converge to the same terminal region (action-cell, goal state, behavioral outcome) when they satisfy the local sufficiency condition.}

At each scale:
\begin{itemize}[noitemsep]
\item \textbf{Molecular}: Multiple enzyme conformations (different internal trajectories) catalyze the same reaction (same terminal state)
\item \textbf{Cellular}: Multiple ion channel states generate the same membrane potential (same terminal electrical state)
\item \textbf{Neural}: Multiple spinal pathways converge to the same motor neuron firing (same terminal firing state)
\item \textbf{Behavioral}: Multiple thought patterns converge to the same action (same behavioral terminus)
\end{itemize}

The nervous system, the body, and cognition are organized around sufficiency, not perfection. This is both the source of robustness (many routes work) and of limitations (many routes are inefficient).
\end{principle}

%=============================================================================
\subsection{Summary: Variational Principle and Convergence}
\label{subsec:variational-summary}

\begin{principle}[Motor Variational Principle]
Motor control can be formalized as minimization of the action functional $\mathcal{S}[\mathbf{u}] = \int \mathcal{L} dt$, where the Lagrangian incorporates:
\begin{itemize}[noitemsep]
\item ATP cost (scaled by local oxygen availability)
\item Mechanical work against external loads
\item Proprioceptive error (expected vs actual feedback)
\item Synaptic state constraints (residual NT and plasticity)
\end{itemize}

Euler-Lagrange equations generate motor command trajectories that extremize this action. Multiple distinct trajectories (spinal reflex, cortical, cerebellar) converge to identical behavioral outcomes when they satisfy the motor sufficiency condition: the gradient of the Lagrangian is small enough to be attracted to the action-cell.

Each re-execution generates a new trajectory (Poincaré deviation $\delta > 0$), creating a historical record. Memory is trajectory history; consciousness is the verification that multiple distinct internal trajectories converge to the same sensory-verified outcome.

This unifies the neuromusculoskeletal system with the fundamental principles of bounded phase space dynamics and categorical convergence established in Sections~\ref{sec:partition-landscape}--\ref{sec:neuromusculoskeletal}.
\end{principle}

